\chapter{Why Compress at All?}
\label{ch:why-compress}

\begin{plainlanguage}
\textbf{In plain terms:} If "intelligence" and "creativity" hide so much structure, why hasn't language already replaced them with more precise vocabulary? The candidate answer: single words survive because they're communicatively \emph{efficient}, not because they're correct. This chapter states that idea precisely and is honest that it doesn't yet prove it.
\end{plainlanguage}

\section{The objection this chapter exists to answer}

A reader who has followed this book to this point has watched intelligence decompose into recognition and implementation, implementation decompose further into a vector of inhibition, emotion regulation, effort, delay, and social cost, and creativity decompose along an overlapping but not identical set of axes. The natural and entirely fair objection at this point is: if these single-word labels are this inaccurate --- if ``intelligence'' and ``creativity'' are each concealing a half-dozen only-partially-correlated sub-competencies --- why do cultures keep using single words for them at all? Why hasn't ordinary language already fragmented into the more accurate vocabulary this book keeps reaching for?

This chapter does not fully answer that question. It states the question precisely enough to be answered, sketches the shape an answer would need to take, and is explicit about stopping short of supplying one. That is the honest status of the material here, and the chapter is better for saying so directly than for dressing an unfinished argument in the language of a completed one.

\section{A candidate account, stated as an open program rather than a result}

\begin{conjecture}[Compression Efficiency]
A competence label persists in a language when, and roughly because, the communicative savings of using one word rather than a full multidimensional profile exceed the predictive error that compression introduces.
\end{conjecture}

Stating this precisely enough to be more than a slogan requires two quantities this chapter can define but not yet specify with any actual content.

\begin{definition}[Description Length and Explanatory Error]
Let $L(n)$ denote the communicative cost --- in whatever unit eventually proves appropriate, plausibly something like processing time, working-memory load, or message length --- of describing a competence using $n$ independent dimensions rather than a single compressed label. Let $E(n)$ denote the predictive or explanatory error that remains when a competence is described using only $n$ dimensions rather than its full underlying structure --- the gap between what an $n$-dimensional description predicts about a person's behavior and what actually happens.
\end{definition}

On this account, a single-word label is communicatively efficient exactly where collapsing from $n$ dimensions to $1$ saves more in $L$ than it costs in $E$:
\[
\Delta L \gg \Delta E \quad \text{(at } n=1\text{, relative to the full decomposition)}.
\]

If something like this inequality holds for ``intelligence'' and ``creativity'' as actually used, it would explain why decomposition keeps happening in careful analysis --- researchers, ethnographers, and this book itself keep finding more structure underneath the label --- while the vocabulary itself remains stable in ordinary use: the label is not false, on this account, so much as it is a genuinely efficient lossy compression, one that most everyday uses of the word do not need to undo, because most everyday predictive purposes (deciding roughly whom to trust with a difficult task, whom to recommend for a job, whom a child should try to emulate) tolerate the error $E(1)$ introduces in exchange for the enormous savings in $L$.

The folk phrases this book has returned to several times --- ``brilliant but self-destructive,'' ``wise but ineffective'' --- are, on this account, exactly the moments where the compression's error becomes locally too costly to ignore, and language responds by reaching for a longer, more decomposed description precisely because the single-word version has stopped being efficient for the purpose at hand. That the longer phrase exists at all, rather than language simply tolerating a wrong prediction, is itself a small piece of evidence that ordinary language is doing something like cost-sensitive compression rather than blind compression.

\subsection*{Why not just keep every distinction?}

A sharper version of this chapter's question is worth asking directly: if compression genuinely costs explanatory accuracy, why does language compress at all, rather than simply maintaining an unbounded, maximally precise vocabulary --- a separate word for every distinguishable combination of recognition and implementation sub-competencies, with nothing ever collapsed into a coarser label?

The candidate answer available from the machinery already built in this chapter is that an infinite vocabulary, even if cognitively achievable, would maximize precision at the direct expense of coordination. Communication requires that a speaker's word and a listener's understanding of that word refer to approximately the same thing without extensive negotiation on each use; the more finely a vocabulary is allowed to fragment, the less likely any two speakers share the exact same fine-grained term for the exact same fine-grained bundle, and the more communicative work --- exactly the cost $L(n)$ this chapter has already introduced --- is required merely to establish shared reference before any actual content can be exchanged. A finite vocabulary trades away some precision in order to purchase a public, shared coordination surface: a word that is somewhat lossy but reliably understood the same way by most speakers most of the time is more useful for the actual business language is for --- coordinating action and belief across many speakers --- than a maximally precise word understood by no one but its coiner.

This connects directly to a theme that recurs elsewhere in work adjacent to this book, on why some distinctions persist across long stretches of time while others are quietly let go: a distinction survives in a shared vocabulary not simply because it is real or because it could in principle be drawn, but because it continues to earn its communicative keep --- because enough speakers, enough of the time, need to coordinate around exactly that boundary for the cost of maintaining a dedicated word for it to be worth paying. An infinite vocabulary would maximize the precision available to a single isolated mind. A finite one is what a vocabulary used by more than one mind, repeatedly, under real time pressure, actually evolves toward.

\section{What is missing, named explicitly}

This is as far as the conjecture can honestly be taken in this book, and what is missing is not a minor gap to be waved past. $L$ and $E$ have been defined in words, not specified with any actual content --- there is no stated unit in which $L$ is measured, no model of what $E$ actually computes, and no derivation showing the inequality $\Delta L \gg \Delta E$ holds for any real competence word rather than merely sounding plausible when described in prose. This is, structurally, the same kind of gap Chapter~\ref{ch:fidelity} left open with the Misweighted Elimination conjecture: a precisely statable, falsifiable-in-principle claim that this book has not closed, offered honestly as the chapter's research target rather than smuggled past as an already-finished result.

This book ends, on this question, not with a theorem but with a research program: a serious, well-motivated conjecture --- that competence words survive not because they are fundamental but because they occupy efficient positions on a communication-cost-versus-explanatory-accuracy frontier --- stated precisely enough that someone, possibly the author of a later book, could build the cost model this chapter does not, and either close the gap or discover that the inequality fails and a different account of why competence words persist is needed instead.

\section{What this completes, and what it leaves open}

Part~V was undertaken to find out whether the machinery built in Parts~II through IV survives contact with a domain it was not built around. Chapter~\ref{ch:creativity} found a genuine, if partial, transport: creativity decomposes along an overlapping but not identical axis to intelligence, supporting the recursive compression conjecture's core claim while complicating its specific two-axis shape. This chapter found that the deeper question the whole program eventually raises --- why compress at all, if compression loses this much --- is real, answerable in principle, and not yet answered here. Both outcomes are reported as what they are. The book attempted its stress test. It did not fully resolve what the stress test surfaced. What survives the attempt --- Intelligence $\neq$ Unique Case, and a precisely stated open conjecture about why competence words persist despite the inaccuracy this book has spent thirteen chapters documenting --- is offered as exactly that much, no more, and the next chapter's resolution of the discovery/constitution fork does not depend on any more than that having been shown.
